\newpage
\section{Analiza problemu}

Rozwój modeli generatywnych spowodował, że obrazy syntetyczne stały się łatwe do wytworzenia i coraz trudniejsze do odróżnienia od zwykłych fotografii. Problem ten ma znaczenie praktyczne w kilku obszarach: dezinformacji, podszywaniu się pod inne osoby, publikowaniu sfałszowanych materiałów wizualnych oraz obniżaniu wiarygodności obrazu jako nośnika informacji. W projekcie przyjęto, że pierwszym krokiem powinno być zbudowanie modelu bazowego wykrywającego obrazy typu \texttt{real} oraz \texttt{fake}.

\subsection{Charakter problemu}

Obrazy generowane przez AI mogą zawierać artefakty trudne do uchwycenia dla człowieka, ale rozpoznawalne przez model uczony na odpowiednio dobranym zbiorze danych. Do takich cech można zaliczyć:

\begin{itemize}
    \item nienaturalną spójność lub niespójność tekstur,
    \item charakterystyczne błędy w drobnych detalach,
    \item zbyt gładkie przejścia tonalne,
    \item artefakty wynikające z procesu generacji lub późniejszej kompresji obrazu.
\end{itemize}

W praktyce problem nie sprowadza się jednak tylko do znalezienia jednego modelu o wysokiej skuteczności. Istotne jest również to, czy model pozostaje odporny na pogorszenie jakości obrazu, takie jak rozmycie, kompresja JPEG czy zmniejszenie rozdzielczości. Z punktu widzenia zastosowań praktycznych właśnie te przypadki są szczególnie ważne, ponieważ obrazy rozpowszechniane w mediach społecznościowych lub komunikatorach rzadko zachowują oryginalną jakość.

\subsection{Cel etapu bazowego}

W pierwszym etapie projektu założono budowę rozwiązania, które:

\begin{itemize}
    \item przyjmuje cały obraz jako wejście,
    \item klasyfikuje go do jednej z dwóch klas: \texttt{real} lub \texttt{fake},
    \item zapisuje pełne metryki i checkpointy,
    \item pozwala na wznowienie treningu,
    \item umożliwia interpretację decyzji modelu za pomocą Grad-CAM,
    \item może zostać później rozszerzone o osobny moduł analizy twarzy.
\end{itemize}

Takie podejście pozwala najpierw zbudować solidny punkt odniesienia, a dopiero później przejść do trudniejszego zagadnienia deepfake twarzy.

\subsection{Najważniejsze trudności}

W trakcie analizy problemu zidentyfikowano kilka kluczowych trudności:

\begin{itemize}
    \item obrazy spoza rozkładu treningowego mogą być klasyfikowane z dużą pewnością, ale błędnie,
    \item portrety niskiej jakości i obrazy po kompresji mogą powodować fałszywe alarmy,
    \item ilustracje i anime nie należą do głównego zakresu modelu fotograficznego,
    \item skuteczność mierzona na zbiorze testowym nie zawsze przekłada się wprost na zachowanie modelu na obrazach z zewnątrz.
\end{itemize}

Z tego względu poza samym treningiem konieczne było dodanie testu odporności oraz prostego mechanizmu ostrzegania o niskiej jakości obrazu w interfejsie demonstracyjnym.
